\documentclass[10pt, aspectratio=169]{beamer}
% Preámbulo modular para presentaciones Beamer (Modelación Estadística)
% Diseño y paleta institucional del Tecnológico de Monterrey

\documentclass[aspectratio=169,xcolor={dvipsnames,table}]{beamer}

\usepackage[utf8]{inputenc}
\usepackage[spanish,mexico]{babel}
\usepackage{amsmath,amssymb,amsfonts,mathtools}
\usepackage{graphicx}
\usepackage{listings}
\usepackage{booktabs}
\usepackage{tikz}
\usetikzlibrary{matrix,arrows,decorations.pathmorphing,shapes.geometric,positioning}

% Paleta Institucional Tecnológico de Monterrey
\definecolor{TecAzul}{HTML}{0039A6}       % Azul TEC: color primario institucional
\definecolor{TecAzulOscuro}{HTML}{1A2E51} % Azul marino oscuro
\definecolor{TecRojo}{HTML}{EC2661}       % Rojo de acento (paleta miTec)
\definecolor{TecGris}{HTML}{646464}       % Gris neutro para comentarios y subtítulos
\definecolor{TecFondoBloque}{HTML}{F4F6F9}% Fondo suave para bloques

% Configuración de Tema Beamer
\usetheme{Boadilla}
\usecolortheme[named=TecAzulOscuro]{structure}
\setbeamercolor{alerted text}{fg=TecRojo}
\setbeamercolor{palette primary}{bg=TecAzulOscuro,fg=white}
\setbeamercolor{palette secondary}{bg=TecAzul,fg=white}
\setbeamercolor{palette tertiary}{bg=TecAzulOscuro!80!black,fg=white}
\setbeamercolor{title}{fg=TecAzulOscuro,bg=TecFondoBloque}
\setbeamercolor{frametitle}{fg=TecAzulOscuro,bg=TecFondoBloque}

% Bloques personalizados elegantes
\setbeamercolor{block title}{bg=TecAzulOscuro,fg=white}
\setbeamercolor{block body}{bg=TecFondoBloque,fg=black}
\setbeamercolor{block title alerted}{bg=TecRojo,fg=white}
\setbeamercolor{block body alerted}{bg=TecRojo!8,fg=black}
\setbeamercolor{block title example}{bg=TecAzul,fg=white}
\setbeamercolor{block body example}{bg=TecAzul!8,fg=black}

% Redondear bordes y añadir sombra sutil
\setbeamertemplate{blocks}[rounded][shadow=true]
\setbeamertemplate{navigation symbols}{}

% Pie de página limpio con número de diapositiva
\setbeamertemplate{footline}{
  \leavevmode%
  \hbox{%
  \begin{beamercolorbox}[wd=.7\paperwidth,ht=2.25ex,dp=1ex,left]{author in head/foot}%
    \hspace*{2ex}\insertshortauthor~~(\insertshortinstitute) \hfill \insertshorttitle
  \end{beamercolorbox}%
  \begin{beamercolorbox}[wd=.3\paperwidth,ht=2.25ex,dp=1ex,right]{date in head/foot}%
    \insertshortdate{}\hspace*{2em}
    \insertframenumber{} / \inserttotalframenumber\hspace*{2ex}
  \end{beamercolorbox}}%
  \vskip0pt%
}

% Configuración de Listings de Python para visualización pedagógica de código
\lstset{
    language=Python,
    basicstyle=\ttfamily\footnotesize,
    keywordstyle=\color{TecAzulOscuro}\bfseries,
    stringstyle=\color{TecRojo},
    commentstyle=\color{TecGris}\itshape,
    showstringspaces=false,
    breaklines=true,
    frame=lines,
    rulecolor=\color{TecAzulOscuro!30},
    backgroundcolor=\color{TecFondoBloque},
    aboveskip=0.5em,
    belowskip=0.5em,
    literate={á}{{\'a}}1 {é}{{\'e}}1 {í}{{\'i}}1 {ó}{{\'o}}1 {ú}{{\'u}}1
             {Á}{{\'A}}1 {É}{{\'E}}1 {Í}{{\'I}}1 {Ó}{{\'O}}1 {Ú}{{\'U}}1
             {ñ}{{\~n}}1 {Ñ}{{\~N}}1 {¿}{{?`}}1 {¡}{{!`}}1
}

% Metadatos institucionales por defecto
\title[Modelación Estadística]{Modelación Estadística para Ciencia de Datos e Ingeniería}
\author[J. Castillo Colmenares]{Juliho David Castillo Colmenares}
\institute[Tec de Monterrey]{Tecnológico de Monterrey \\ Escuela de Ingeniería y Ciencias}
\date{\today}

% Comandos y macros estadísticas compartidas para las presentaciones Beamer (Metropolis)

% Conjuntos numéricos básicos
\newcommand{\R}{\mathbb{R}}
\newcommand{\N}{\mathbb{N}}
\newcommand{\Q}{\mathbb{Q}}
\newcommand{\Z}{\mathbb{Z}}
\newcommand{\set}[1]{\left\{ #1 \right\}}
\newcommand{\abs}[1]{\left|#1\right|}
\newcommand{\norm}[1]{\left\| #1 \right\|}

% Comandos estadísticos y probabilísticos del curso
\newcommand{\Var}{\operatorname{Var}}
\newcommand{\cov}{\operatorname{Cov}}
\newcommand{\corr}{\rho}
\newcommand{\comb}[2]{\begin{pmatrix} #1 \\ #2 \end{pmatrix}}
\newcommand{\s}{\sigma}
\newcommand{\particion}{\mathfrak{P}}
\newcommand{\card}[1]{\#\left( #1 \right)}
\newcommand{\seq}[3]{\set{#1}_{#2}^{#3}}
\newcommand{\E}{\mathbb{E}}
\newcommand{\Prob}{\mathbb{P}}

% Entornos pedagógicos y cajas en español académico natural
\newenvironment{discusion}{%
  \begin{alertblock}{Pregunta de Discusión}%
}{%
  \end{alertblock}%
}

\newenvironment{reflexion}{%
  \begin{alertblock}{Reflexión para el Aula}%
}{%
  \end{alertblock}%
}

\newenvironment{paradoja}{%
  \begin{alertblock}{Paradoja de Exploración}%
}{%
  \end{alertblock}%
}

\newenvironment{motivacion}{%
  \begin{exampleblock}{Motivación: Ciencia de Datos en la Práctica}%
}{%
  \end{exampleblock}%
}

\newenvironment{retoclase}{%
  \begin{block}{Reto Rápido en Equipo}%
}{%
  \end{block}%
}


\title[Geométrica y Binomial Negativa]{Sección 03.05: Geométrica y Binomial Negativa}
\subtitle{Tiempos de Espera, Pérdida de Memoria y Sobredispersión}
\author[J. Castillo Colmenares]{Juliho Castillo Colmenares}
\institute{Tecnológico de Monterrey}
\date{\vspace{-2.0cm}}

\begin{document}

% Slide 1: Portada
\begin{frame}[plain]
  \vspace{-0.5cm}
  \titlepage
\end{frame}

% Slide 2: Hoja de Ruta Modular
\begin{frame}{Hoja de Ruta — Capítulo 03: Variables Aleatorias Discretas}
  \footnotesize
  ¿Dónde nos encontramos en el desarrollo modular de la teoría? \pause
  \begin{itemize}\itemsep=0.08em
    \item \textbf{03.01} Concepto, Notación y Definición de Variable Aleatoria \pause
    \item \textbf{03.02} Funciones de Probabilidad (PMF) y Distribución Acumulada (CDF) \pause
    \item \textbf{03.03} Esperanza Matemática y Varianza Operativa \pause
    \item \textbf{03.04} Distribución Binomial y Ensayos de Bernoulli \pause
    \item \color{TecRojo}\textbf{03.05 Distribuciones Geométrica y Binomial Negativa \emph{(Hoy)}}\color{black} \pause
    \item \textbf{03.06} Distribución Hipergeométrica \pause
    \item \textbf{03.07} Distribución de Poisson y Procesos de Poisson
  \end{itemize}
  \vspace{-0.1cm}
  \begin{block}{Objetivo de la Sesión}
    Analizar el modelado de tiempos de espera hasta el $r$-ésimo éxito en ensayos de Bernoulli, dominando la propiedad de pérdida de memoria para $r=1$ y el fenómeno de sobredispersión ($\Var(X) > \E[X]$) para $r \ge 1$.
  \end{block}
\end{frame}

% Slide 3: Motivación en Ciencia de Datos
\begin{frame}{Motivación — De Contar Éxitos a Medir Tiempos de Espera}
  \footnotesize
  \begin{motivacion}[¿Por qué cambiar la perspectiva de la variable aleatoria?]
    En ingeniería de confiabilidad o ciberseguridad, no siempre nos interesa saber cuántas fallas habrá en un número fijo de horas (Binomial). Con frecuencia, la pregunta crítica para el negocio es: \emph{¿Cuántas transacciones u horas debemos esperar hasta que ocurra el primer ataque o la cuarta falla crítica?}
  \end{motivacion}
  
  \vspace{0.15cm}
  \pause
  \begin{columns}[T]
    \begin{column}{0.48\textwidth}
      \begin{block}{Enfoque Binomial (Ensayos Fijos)}
        \scriptsize
        \begin{itemize}\itemsep=0.06cm
          \item Fijamos el total de intentos $n$.
          \item $X = \text{éxitos observados}$ ($X \le n$).
          \item Soporte finito: $\{0, 1, \dots, n\}$.
        \end{itemize}
      \end{block}
    \end{column}
    \begin{column}{0.48\textwidth}
      \begin{alertblock}{Enfoque de Espera (Éxitos Fijos)}
        \scriptsize
        \begin{itemize}\itemsep=0.06cm
          \item Fijamos la meta de éxitos $r$.
          \item $X = \text{ensayos necesarios}$ ($X \ge r$).
          \item Soporte infinito numerable: $\{r, r+1, \dots\}$.
        \end{itemize}
      \end{alertblock}
    \end{column}
  \end{columns}
\end{frame}

% Slide 4: Distribución Geométrica
\begin{frame}{La Distribución Geométrica ($X \sim \text{Geom}(p)$)}
  \footnotesize
  Sea $X$ el número total de ensayos de Bernoulli independientes e idénticamente distribuidos (i.i.d.) requeridos para obtener el \textbf{primer éxito} con probabilidad $p \in (0, 1)$ y fracaso $q = 1-p$: \pause

  \vspace{0.1cm}
  \begin{block}{Función de Masa de Probabilidad (PMF) y Soporte}
    $$P(X = k) = (1-p)^{k-1}p = q^{k-1}p, \quad \text{para } k \in \mathcal{S}_X = \{1, 2, 3, \dots\}$$
  \end{block}

  \vspace{0.1cm}
  \pause
  \begin{alertblock}{Normalización Exacta y Función Acumulada (CDF)}
    \begin{itemize}\itemsep=0.06cm
      \item \textbf{Normalización:} $\sum_{k=1}^{\infty} q^{k-1}p = p \sum_{j=0}^{\infty} q^j = p \left(\frac{1}{1-q}\right) = \frac{p}{p} = 1$. \pause
      \item \textbf{Cola Superior:} $P(X > k) = q^k$ (la probabilidad de fallar los primeros $k$ intentos).
      \item \textbf{CDF:} $F_X(k) = P(X \le k) = 1 - q^k$. En Python (`scipy.stats`): \texttt{geom.cdf(k, p)}.
    \end{itemize}
  \end{alertblock}
\end{frame}

% Slide 5: Momentos Geométricos y Pérdida de Memoria
\begin{frame}{Momentos y Propiedad de Pérdida de Memoria}
  \footnotesize
  A partir de la función generadora de momentos $M_X(t) = \frac{pe^t}{1-qe^t}$, obtenemos los momentos descriptivos de la ley geométrica: \pause

  \vspace{0.1cm}
  \begin{block}{Esperanza y Varianza Operacional}
    $$\E[X] = \dfrac{1}{p}, \qquad \Var(X) = \dfrac{q}{p^2} = \dfrac{1-p}{p^2}$$
  \end{block}

  \vspace{0.1cm}
  \pause
  \begin{alertblock}{Propiedad de Pérdida de Memoria (Markoviana Discreta)}
    La distribución geométrica es la \textbf{única ley discreta sobre $\Z^+$} sin memoria respecto al pasado:
    $$P(X > m+n \mid X > m) = \dfrac{P(X > m+n)}{P(X > m)} = \dfrac{q^{m+n}}{q^m} = q^n = P(X > n)$$
    \scriptsize\emph{Interpretación:} Haber esperado $m$ intentos fallidos no cambia la distribución del tiempo restante.
  \end{alertblock}
\end{frame}

% Slide 6: Distribución Binomial Negativa
\begin{frame}{La Distribución Binomial Negativa ($X \sim \text{BN}(r, p)$)}
  \footnotesize
  Generalicemos el tiempo de espera: sea $X$ el número total de ensayos requeridos para acumular \textbf{exactamente $r$ éxitos} independientes ($r \ge 1$): \pause

  \vspace{0.1cm}
  \begin{block}{Función de Masa de Probabilidad (Ley de Pascal)}
    $$P(X = k) = \comb{k-1}{r-1} p^r (1-p)^{k-r}, \quad \text{para } k \in \mathcal{S}_X = \{r, r+1, r+2, \dots\}$$
  \end{block}

  \vspace{0.1cm}
  \pause
  \begin{alertblock}{Descomposición Combinatoria y Estocástica}
    \begin{itemize}\itemsep=0.06cm
      \item \textbf{Lógica Combinatoria:} Para que el intento $k$ sea el $r$-ésimo éxito, debe haber exactly $r-1$ éxitos distribuidos en los $k-1$ ensayos previos ($\comb{k-1}{r-1}p^{r-1}q^{k-r}$) multiplicado por el éxito final ($p$).
      \item \textbf{Suma Estocástica:} $X$ equivale exactamente a la suma de $r$ variables geométricas independientes: $X = \sum_{i=1}^r Y_i$, donde $Y_i \sim \text{Geom}(p)$.
    \end{itemize}
  \end{alertblock}
\end{frame}

% Slide 7: Momentos y Sobredispersión
\begin{frame}{Momentos de la Binomial Negativa y Sobredispersión}
  \footnotesize
  Por linealidad de la esperanza y aditividad de varianzas para variables independientes ($Y_i \sim \text{Geom}(p)$): \pause

  \vspace{0.1cm}
  \begin{block}{Esperanza y Varianza del Tiempo Total}
    $$\E[X] = \sum_{i=1}^r \E[Y_i] = \dfrac{r}{p}, \qquad \Var(X) = \sum_{i=1}^r \Var(Y_i) = \dfrac{rq}{p^2} = \dfrac{r(1-p)}{p^2}$$
  \end{block}

  \vspace{0.1cm}
  \pause
  \begin{alertblock}{Análisis de Sobredispersión ($\Var(X) > \E[X]$)}
    En cualquier variable discreta, el ratio de dispersión $\Var(X) / \E[X]$ caracteriza la estructura del modelo:
    \begin{itemize}\itemsep=0.04cm
      \item \textbf{Binomial (`Sub-dispersa`):} $\Var(X) / \E[X] = 1-p < 1$.
      \item \textbf{Poisson (`Equi-dispersa`):} $\Var(X) / \E[X] = 1$.
      \item \textbf{Binomial Negativa (`Sobre-dispersa`):} $\Var(X) / \E[X] = \frac{1}{p} > 1$. Ideal para datos con alta variabilidad.
    \end{itemize}
  \end{alertblock}
\end{frame}

% Slide 8: Relación Asintótica Pascal a Poisson
\begin{frame}{Relación Asintótica — De Pascal a Poisson}
  \footnotesize
  Consideremos la parametrización del número de fracasos observados antes de los $r$ éxitos: $Y = X - r \sim \text{BN}(r, p)$, con media $\mu_Y = \frac{rq}{p}$. \pause

  \vspace{0.1cm}
  \begin{block}{Teorema del Límite de Eventos Raros ($r \to \infty, \, p \to 1$)}
    Si la exigencia de éxitos crece indefinidamente ($r \to \infty$) pero cada intento es casi seguro ($p \to 1$), manteniendo la media de fracasos constante e igual a una tasa fija $(1-p)r = \lambda$:
    $$\lim_{r \to \infty} \comb{y+r-1}{y} p^r (1-p)^y = \dfrac{\lambda^y e^{-\lambda}}{y!}$$
  \end{block}

  \vspace{0.1cm}
  \pause
  \begin{alertblock}{Relevancia en Modelación Continua}
    Cuando un proceso presenta sobredispersión moderada, la Binomial Negativa actúa como una \textbf{versión robusta y generalizada de la ley de Poisson}, permitiendo capturar heterogeneidad entre individuos.
  \end{alertblock}
\end{frame}

% Slide 9A: Lab Python I (Setup y PMF/CDF)
\begin{frame}[fragile]{Laboratorio en Python: Validación Exacta PMF y CDF}
  \scriptsize
  Verificación en SciPy (\texttt{03.05\_geometric\_negative\_binomial.py}) de la exactitud vectorizada de la PMF y CDF:
  \vspace{0.1cm}
  {\lstset{basicstyle=\fontsize{6pt}{7pt}\selectfont\ttfamily}
  \lstinputlisting[firstline=16, lastline=37]{../../code/03_variables_aleatorias_discretas/03.05_geometric_negative_binomial.py}}
\end{frame}

% Slide 9B: Lab Python II (Monte Carlo Pérdida de Memoria)
\begin{frame}[fragile]{Laboratorio en Python: Pérdida de Memoria por Monte Carlo}
  \scriptsize
  Simulación de $250,000$ ensayos verificando empíricamente la propiedad sin memoria $P(X > m+n \mid X > m) = P(X > n)$:
  \vspace{0.08cm}
  {\lstset{basicstyle=\fontsize{6pt}{7pt}\selectfont\ttfamily}
  \lstinputlisting[firstline=42, lastline=58]{../../code/03_variables_aleatorias_discretas/03.05_geometric_negative_binomial.py}}
\end{frame}

% Slide 9C: Lab Python III (Sobredispersión y Aditividad)
\begin{frame}[fragile]{Laboratorio en Python: Sobredispersión y Aditividad}
  \scriptsize
  Demostración empírica de la aditividad geométrica y del ratio de sobredispersión estrictamente mayor a 1:
  \vspace{0.08cm}
  {\lstset{basicstyle=\fontsize{6pt}{7pt}\selectfont\ttfamily}
  \lstinputlisting[firstline=63, lastline=77]{../../code/03_variables_aleatorias_discretas/03.05_geometric_negative_binomial.py}}
\end{frame}

% Slide 9D: Lab Python IV (Salida Consola)
\begin{frame}[fragile]{Laboratorio en Python: Salida en Terminal}
  \scriptsize
  Salida estándar generada al ejecutar el script del laboratorio en Python:
  \vspace{0.05cm}
  \begin{block}{Salida Estándar en Consola}
    \fontsize{6pt}{7pt}\selectfont
    \begin{verbatim}
=== Section 03.05: Computational Lab Verification ===
[Snippet 1] Max Geometric PMF Error against SciPy: 0.00e+00
[Snippet 1] Max Negative Binomial PMF Error against SciPy: 1.39e-17

[Snippet 2] Memoryless Verification (p=0.25, m=5, n=3):
  Empirical Conditional P(X > 8 | X > 5): 0.420767
  Empirical Marginal P(X > 3):            0.422532
  Exact Theoretical Probability (1-p)^3:  0.421875

[Snippet 3] Overdispersion & Additivity Verification (p=0.20, r=5):
  Empirical Mean: 25.0140 | Theoretical Mean: 25.0000
  Empirical Var:  100.4530 | Theoretical Var:  100.0000
  Overdispersion Ratio (Var / Mean): 4.0159 (> 1.0 confirms overdispersion)

SUCCESS: All computational checks passed cleanly.
    \end{verbatim}
  \end{block}
\end{frame}

% Slide 10A: Ejercicio en Clase — Nivel 1: Fundamental (Enunciado)
\begin{frame}{Ejercicio en Clase — Nivel 1: Fundamental (1/2)}
  \footnotesize
  \begin{block}{Problema 3.5.2 — Servidor de Base de Datos y Pérdida de Memoria ($r=1$)}
    En un servidor de base de datos de alta concurrencia, las transacciones independientes experimentan fallas de bloqueo con probabilidad constante $p = 0.04$ ($q = 0.96$). Sea $X$ el número de intentos necesarios para que una transacción logre ejecutarse con éxito por primera vez.
  \end{block}

  \vspace{0.2cm}
  \pause
  \begin{alertblock}{Planteamiento de Incisos}
    \begin{enumerate}\itemsep=0.1cm
      \item Calcule la probabilidad exacta de que la transacción requiera a lo sumo $3$ intentos para tener éxito ($P(X \le 3)$).
      \item Si la transacción ya ha fallado sus primeros $5$ intentos, ¿cuál es la probabilidad condicional de que requiera más de $8$ intentos en total ($P(X > 8 \mid X > 5)$)?
      \item Compárela con $P(X > 3)$ e interprete el resultado.
    \end{enumerate}
  \end{alertblock}
\end{frame}

% Slide 10B: Ejercicio en Clase — Nivel 1: Fundamental (Resolución)
\begin{frame}{Ejercicio en Clase — Nivel 1: Fundamental (2/2)}
  \scriptsize
  Dado que $X \sim \text{Geom}(p=0.04)$, utilizamos las fórmulas exactas de la CDF y cola superior: \pause

  \vspace{0.04cm}
  \begin{block}{Evaluación de Probabilidades Acumuladas}
    \begin{itemize}\itemsep=0.04cm
      \item \textbf{A lo sumo 3 intentos ($X \le 3$):}
      $$P(X \le 3) = F_X(3) = 1 - q^3 = 1 - (0.96)^3 = 1 - 0.884736 = \mathbf{0.115264}$$ \pause
      \item \textbf{Probabilidad condicional tras 5 fallas:} Por definición o por propiedad sin memoria ($m=5, n=3$):
      $$P(X > 5+3 \mid X > 5) = P(X > 3) = q^3 = (0.96)^3 = \mathbf{0.884736}$$
    \end{itemize}
  \end{block}

  \vspace{0.04cm}
  \pause
  \begin{alertblock}{Conclusión de Confiabilidad}
    El servidor no tiene memoria de las fallas pasadas: el hecho de haber fallado 5 veces no incrementa ni disminuye la probabilidad de éxito en el sexto intento, la cual permanece exactamente en $4\%$.
  \end{alertblock}
\end{frame}

% Slide 11A: Ejercicio en Clase — Nivel 2: Operativo (Enunciado)
\begin{frame}{Ejercicio en Clase — Nivel 2: Operativo (1/2)}
  \footnotesize
  \begin{block}{Problema 3.5.6 — Tráfico de Red y Sobredispersión por Método de Momentos ($r>1$)}
    Un ingeniero analiza reintentos de conexión en hora punta. Al ajustar el número de fallas previas al primer éxito por usuario ($Y = X - r \sim \text{BN}(r, p)$), observa una media muestral $\bar{x} = 3.2$ fallas y una varianza muestral $s^2 = 13.44$.
  \end{block}

  \vspace{0.2cm}
  \pause
  \begin{alertblock}{Incisos del Diagnóstico de Red}
    \begin{enumerate}\itemsep=0.1cm
      \item Explique teóricamente por qué un modelo de Poisson pura ($\mu = \sigma^2$) resulta estructuralmente inadecuado para este tráfico.
      \item Sabiendo que para $Y \sim \text{BN}(r, p)$ se tiene $\mu_Y = \frac{rq}{p}$ y $\sigma_Y^2 = \frac{rq}{p^2}$, utilice el método de momentos igualando con $\bar{x}$ y $s^2$ para estimar los parámetros poblacionales $\hat{p}$ y $\hat{r}$.
    \end{enumerate}
  \end{alertblock}
\end{frame}

% Slide 11B: Ejercicio en Clase — Nivel 2: Operativo (Resolución)
\begin{frame}{Ejercicio en Clase — Nivel 2: Operativo (2/2)}
  \fontsize{7pt}{8.5pt}\selectfont
  Analizamos el cociente muestral de dispersión para diagnosticar el tráfico: \pause

  \vspace{-0.05cm}
  \begin{block}{Diagnóstico de Dispersión y Estimación de Parámetros}
    \fontsize{7pt}{8.5pt}\selectfont
    \begin{itemize}\itemsep=0.08cm
      \item \textbf{Inviabilidad de Poisson:} El ratio muestral es $s^2 / \bar{x} = 13.44 / 3.2 = \mathbf{4.20} \gg 1$. Existe una fuerte sobredispersión que la ley de Poisson jamás podría reproducir. \pause
      \item \textbf{Estimación de $\hat{p}$ (Cociente de Momentos):}
      $\dfrac{\sigma_Y^2}{\mu_Y} = \dfrac{rq/p^2}{rq/p} = \dfrac{1}{p} \implies \hat{p} = \dfrac{\bar{x}}{s^2} = \dfrac{3.2}{13.44} \approx \mathbf{0.23809}$ \pause
      \item \textbf{Estimación de $\hat{r}$ (Meta de Éxitos):}
      $\mu_Y = \dfrac{r(1-p)}{p} \implies \hat{r} = \dfrac{\bar{x} \cdot \hat{p}}{1 - \hat{p}} = \dfrac{3.2(0.23809)}{0.76190} = \mathbf{1.0000}$
    \end{itemize}
  \end{block}

  \vspace{-0.1cm}
  \pause
  \begin{alertblock}{Conclusión de Telecomunicaciones}
    \fontsize{7pt}{8.5pt}\selectfont
    El tráfico se modela perfectamente con $r=1$ y $p=0.238$, revelando que la alta varianza es consecuencia natural de una distribución geométrica latente de reintentos.
  \end{alertblock}
\end{frame}

% Slide 12A: Ejercicio en Clase — Nivel 3: Analítico (Enunciado)
\begin{frame}{Ejercicio en Clase — Nivel 3: Analítico (1/2)}
  \footnotesize
  \begin{block}{Problema 3.5.7 — Descomposición Estocástica y Función Generadora de Momentos}
    Sea un experimento secuencial donde $X \sim \text{BN}(r, p)$ cuenta el número total de ensayos de Bernoulli necesarios para acumular $r$ éxitos independientes ($r \ge 1$).
  \end{block}

  \vspace{0.2cm}
  \pause
  \begin{alertblock}{Incisos de Demostración Analítica}
    \begin{enumerate}\itemsep=0.1cm
      \item Demuestre por descomposición estocástica que $X$ puede expresarse como la suma de $r$ variables aleatorias geométricas independientes ($X = \sum_{i=1}^r Y_i$, donde $Y_i \sim \text{Geom}(p)$).
      \item Sabiendo que para una variable geométrica $M_{Y_1}(t) = \frac{p e^t}{1 - q e^t}$ para $t < -\ln(q)$, deduzca utilizando independencia la función generadora de momentos $M_X(t)$ de la Binomial Negativa.
    \end{enumerate}
  \end{alertblock}
\end{frame}

% Slide 12B: Ejercicio en Clase — Nivel 3: Analítico (Resolución)
\begin{frame}{Ejercicio en Clase — Nivel 3: Analítico (2/2)}
  \scriptsize
  Demostramos formalmente desglosando la trayectoria secuencial del proceso: \pause

  \vspace{-0.1cm}
  \begin{block}{Paso 1: Tiempos Entre Éxitos Sucesivos}
    \scriptsize
    Sea $Y_1$ los intentos hasta el 1er éxito, y $Y_i$ los intentos desde el éxito $i-1$ hasta el éxito $i$. Por independencia de los ensayos de Bernoulli, cada $Y_i$ es i.i.d. y sigue una ley $\text{Geom}(p)$. El tiempo total para alcanzar $r$ éxitos es la suma exacta de sus esperas parciales: $X = \sum_{i=1}^r Y_i$.
  \end{block}

  \vspace{-0.15cm}
  \pause
  \begin{alertblock}{Paso 2: Producto de Funciones Generadoras de Momentos}
    \scriptsize
    Por el teorema central de variables aleatorias independientes, la MGF de una suma es el producto de las MGFs individuales:
    $$M_X(t) = \E[e^{t X}] = \E\left[ \prod_{i=1}^r e^{t Y_i} \right] = \prod_{i=1}^r \E[e^{t Y_i}] = \prod_{i=1}^r M_{Y_i}(t) = \left( \dfrac{p e^t}{1 - q e^t} \right)^r \quad \blacksquare$$
  \end{alertblock}
\end{frame}

% Slide 13A: Ejercicio en Clase — Nivel 4: Desafiante (Enunciado)
\begin{frame}{Ejercicio en Clase — Nivel 4: Desafiante (1/2)}
  \footnotesize
  \begin{block}{Problema 3.5.10 — Límite Asintótico de Pascal a Poisson}
    Considere la parametrización del número de fracasos antes de $r$ éxitos: $Y = X - r \sim \text{BN}(r, p)$, con $P(Y = y) = \comb{y+r-1}{y} p^r q^y$. Suponga que la exigencia de éxitos crece ($r \to \infty$) mientras la probabilidad individual tiende a la certeza ($p \to 1$), manteniendo la media de fracasos fija y constante $(1-p)r = \lambda > 0 \implies q = \frac{\lambda}{r}$.
  \end{block}

  \vspace{0.2cm}
  \pause
  \begin{alertblock}{Objetivo de la Demostración}
    Demuestre analíticamente mediante límites combinatorios y exponenciales que para cualquier entero no negativo $y \in \{0, 1, \dots\}$:
    $$\lim_{r \to \infty} P(Y = y) = \dfrac{\lambda^y e^{-\lambda}}{y!}$$
  \end{alertblock}
\end{frame}

% Slide 13B: Ejercicio en Clase — Nivel 4: Desafiante (Resolución)
\begin{frame}{Ejercicio en Clase — Nivel 4: Desafiante (2/2)}
  \fontsize{7pt}{8.5pt}\selectfont
  Sustituimos $p = 1 - \frac{\lambda}{r}$ y $q = \frac{\lambda}{r}$ en la PMF, y expandimos el coeficiente combinatorio: \pause

  \vspace{-0.15cm}
  \begin{block}{Paso 1: Expansión Algebraica}
    \fontsize{7pt}{8.5pt}\selectfont
    $$P(Y = y) = \dfrac{(r+y-1)(r+y-2)\cdots r}{y!} \left(1 - \dfrac{\lambda}{r}\right)^r \left(\dfrac{\lambda}{r}\right)^y = \dfrac{\lambda^y}{y!} \left(1 - \dfrac{\lambda}{r}\right)^r \prod_{j=0}^{y-1} \left(1 + \dfrac{j}{r}\right)$$
  \end{block}

  \vspace{-0.15cm}
  \pause
  \begin{alertblock}{Paso 2: Evaluación del Límite Simultáneo}
    \fontsize{7pt}{8.5pt}\selectfont
    Tomando el límite cuando $r \to \infty$ para un entero $y$ finito fijo:
    \begin{itemize}\itemsep=0.01cm
      \item Producto combinatorio: $\lim_{r \to \infty} \prod_{j=0}^{y-1} (1 + \frac{j}{r}) = 1 \cdot 1 \cdots 1 = 1$.
      \item Límite exponencial de Euler: $\lim_{r \to \infty} (1 - \frac{\lambda}{r})^r = e^{-\lambda} \implies \lim_{r \to \infty} P(Y=y) = \dfrac{\lambda^y e^{-\lambda}}{y!} \quad \blacksquare$
    \end{itemize}
  \end{alertblock}
\end{frame}

% Slide 14: Cierre de Sección
\begin{frame}{Cierre de Sección y Síntesis Estocástica}
  \begin{reflexion}[El Criterio de Selección de Modelos Discretos]
    El estudio de la familia Geométrica y Binomial Negativa nos brinda la herramienta definitiva para modelar la incertidumbre en procesos de conteo e intervalos de espera. Mientras que la Binomial impone un techo duro en el soporte ($n$) y sub-dispersión, la Binomial Negativa abre el soporte a $\Z^+$ y captura la heterogeneidad del mundo real mediante la sobredispersión.
  \end{reflexion}

  \vspace{0.3cm}
  \pause
  \begin{block}{Resumen de Propiedades Fundamentales}
    \begin{itemize}\itemsep=0.12cm
      \item \textbf{Pérdida de Memoria ($r=1$):} El pasado no condiciona la probabilidad del futuro inmediato.
      \item \textbf{Aditividad Estocástica:} $\text{BN}(r, p)$ es la suma de $r$ variables geométricas independientes.
      \item \textbf{Puente a Poisson:} Cuando $r \to \infty$ y $p \to 1$, los tiempos de espera macroscópicos convergen a una tasa de eventos raros $\lambda$.
    \end{itemize}
  \end{block}
\end{frame}

% Slide 15: Reto para la Próxima Sesión
\begin{frame}{Perspectiva Modular — El Próximo Reto}
  \scriptsize
  \begin{retoclase}[Para la próxima sesión: Distribución Hipergeométrica]
    Hasta este punto, todos los modelos estudiados (Bernoulli, Binomial, Geométrica, Binomial Negativa) asumen rigurosamente que los ensayos son \textbf{mutuamente independientes} y que la probabilidad de éxito $p$ permanece inmutable (muestreo \emph{con reemplazo} o en población infinita). ¿Qué sucede cuando extraemos elementos \emph{sin reemplazo} de un lote cerrado de tamaño finito $N$?
  \end{retoclase}

  \vspace{0.08cm}
  \pause
  \begin{alertblock}{Preparación Recomendada}
    \begin{itemize}\itemsep=0.06cm
      \item Reflexiona sobre cómo cambia la probabilidad en una segunda extracción si la primera fue un éxito ($P(E_2 \mid E_1)$ vs $P(E_2 \mid E_1^c)$).
      \item Repasa las fórmulas combinatorias básicas de selección sin orden ($\comb{N}{n}$) que forman la base del muestreo exacto en auditoría de calidad.
    \end{itemize}
  \end{alertblock}
\end{frame}

\end{document}
